\chapter{Introduction}
\label{chap:introduction}

\section{Motivation}

Large language models increasingly operate as agents: they observe an
environment, choose tools, inspect the result, and continue until a task is
complete. Browser agents apply this loop to navigation, form entry, file
handling, and web applications in which a model output can become an external
side effect. This capability is useful precisely because web tasks are rarely
one-shot questions: benchmarks such as WebArena require agents to carry state
across several dependent interactions~\cite{zhou2024webarena}.

The same loop creates a security problem. Web observations mix application
content, user-generated data, embedded resources, and text controlled by an
attacker. An indirect prompt injection places instructions in that data so the
agent encounters them while pursuing a legitimate goal~\cite{greshake2023indirect}.
The attack exploits an authority ambiguity rather than a malformed API call:
the model is expected to understand novel page content, yet must also decide
whether that content is entitled to redirect its behaviour. Once the model can
invoke privileged tools, a successful redirection can change remote state or
disclose information instead of merely producing incorrect prose.
Prompt reminders, delimiters, and detectors can reduce exposure, but they
leave the acting model responsible for respecting the distinction between the
operator's goal and hostile data. Classical protection principles suggest a
different allocation of responsibility. Complete mediation places a check at
each privileged use, while least privilege limits which component may perform
the effect~\cite{saltzer1975protection}. For an agent, the model can remain a
flexible planner without being the final authority for its proposed actions.
An independent runtime can enforce deterministic constraints, control data
flows, request narrowly scoped authority, and record the outcome.

\section{Research question and thesis}

The research question is:

\begin{quote}
\textbf{RQ.} \rqone
\end{quote}

Security and utility are evaluated together. A system that blocks every action
can suppress attacks while destroying the purpose of an agent; conversely,
high task completion is not useful if trajectories also achieve attacker
goals. The dissertation therefore measures native user-task success, native
attacker success, their joint outcome, and whether attacker-directed actions
reach the executor.

The thesis is that \emph{independent runtime enforcement can reduce evaluated
policy violations while preserving useful autonomy under the tested
conditions}. Independence means that an allow, block, approval, or takeover
outcome is determined outside the acting model. The claim is tied to defined
tool and browser benchmarks rather than to prompt injection in general.

\section{Approach}

This dissertation presents \systemname, a local research prototype positioned
between a BrowserUse agent and its privileged executor. It governs the agent
loop at two points. Before inference, an observation guard identifies
untrusted browser, extraction, and tool-result regions; detects configured
prompt-injection signals; removes attributable regions; and re-detects the
exact reconstructed model input. Before execution, an action guard rebuilds
live browser context, classifies capabilities and sensitive data, propagates
origin-scoped taint, and invokes one policy engine. The executor is reached
only after that decision.

Automatic enforcement is combined with selective human authority. Hard policy
and protected data flows remain non-overridable. A permissible but
consequential action can receive narrow, expiring approval and is evaluated
again against live state before execution. Recognised credentials, one-time
codes, payment data, and other protected values trigger direct takeover so the
operator performs the step instead of authorising the agent to type the value.
If a suspicious current-page observation cannot be sanitised and verified, it
is withheld and the operator may inspect the page before Sentinel captures a
fresh observation. A Chrome extension presents these states and the redacted
decision history; the server remains the source of authority.

Evaluation proceeds in two stages. AgentDojo supplies controlled tool-agent
tasks with native user and attacker oracles across four application domains
~\cite{debenedetti2024agentdojo}. WASP then exercises a BrowserUse agent over
mutable GitLab and Reddit environments, adding observation construction,
navigation, and browser re-planning~\cite{evtimov2025wasp}. The primary browser
result is a matched Qwen 3.5 27B off/enforce comparison. A provenance-separated
9B development sequence is retained because it explains why warning-only
enforcement and incomplete sanitisation were replaced by post-verified
removal.

\section{Contributions}
\label{sec:contributions}

The dissertation makes three contributions.

\begin{enumerate}
  \item \textbf{Executor-side enforcement for a dynamic browser-agent loop.}
  Sentinel interposes on model input and every model-proposed privileged action
  entering the governed executor. Live context, central policy evaluation,
  fail-closed outcomes, and paired attempted/terminal audit separate a model's
  proposal from authority to act without changing the underlying model.

  \item \textbf{Post-verified observation governance with selective
  authority.} Detection is connected to attributable sanitisation,
  reconstruction, and clean re-detection. Origin-scoped taint and action policy
  remain as backstops. Approval, direct takeover, and unresolved-observation
  incidents provide distinct server-owned paths for decisions that cannot be
  resolved as an automatic allow or block.

  \item \textbf{A two-environment security--utility evaluation.} Across 629
  attacked AgentDojo tasks per condition and model, enforcement reduced native
  attacker success from 206 to 0 cases for Qwen 3.5 4B, from 131 to 1 for 9B,
  and from 203 to 0 for 27B while utility increased in each row. In the matched
  27B WASP comparison, attacker success changed from 5/84 to 0/84 and
  attacker-directed execution from 17/84 to 0/84, while utility changed from
  47/84 to 46/84. Reporting intermediate execution together with utility avoids
  treating either inactivity or baseline model failure as security.
\end{enumerate}

Sentinel does not claim BrowserUse, AgentDojo, WASP, or Qwen as original work.
Its original contribution is the architecture that composes observation
governance, executor-side policy, selective authority, and paired audit around
a dynamic browser agent, together with the implementation and evaluation of
that composition.

\section{Dissertation structure}

\Cref{chap:background} establishes the prompt-injection and system-security
context. \Cref{chap:threat-model} defines the authority boundary and security
requirements. \Cref{chap:design} presents the implementation and operator
control plane. \Cref{chap:methodology,chap:results} define and report the
evaluation, and \cref{chap:discussion,chap:conclusion} interpret the evidence,
scope its validity, and state the resulting contribution.
